\section{A controlled dual-domain Push-T testbed}
\label{sec:platform}

Testing the theory requires a testbed in which every component of the embodiment gap is independently manipulable --- something real egocentric datasets do not allow.
We therefore build a \emph{dual-domain} variant of Push-T, following the simulation-side design of EgoBridge~\cite{egobridge2025} with three deliberate deviations (stated below), on top of the diffusion-policy codebase of~\cite{chi2023diffusion}.

\paragraph{Domain construction.}
Both domains share the Push-T task (push a T-shaped block to a goal pose with a point-mass pusher, image observations) and differ along four axes, summarized in Table~\ref{tab:domains} and Fig.~\ref{fig:teaser}: the pusher embodiment (triangle vs.\ circle), the appearance (white vs.\ purple background), the contact dynamics (friction scaled $\times 0.7$ vs.\ unchanged), and the task geometry (the T-block's initial and goal poses are mirrored about the vertical axis in the human domain, so that optimal pushing actions are \emph{not} directly transferable --- a controlled kinematic gap).
The four axes map onto the model's quantities: appearance/friction differences reshape the two spectra $(a_j, b_j)$; the mirrored geometry induces the task mismatch $\delta$; the pusher change alters both.

\begin{table}[t]
  \centering
  \small
  \caption{Dual-domain Push-T construction. The robot domain is the training+evaluation target; the human domain is the mixture source (a \emph{surrogate} for egocentric data, in the spirit of~\cite{egobridge2025}).}
  \label{tab:domains}
  \begin{tabular}{lcc}
    \toprule
     & robot (target) & human (source) \\
    \midrule
    pusher & triangle & circle \\
    background & white & purple \\
    friction & $\times 0.7$ & $\times 1.0$ \\
    T-block geometry & canonical & mirrored \\
    demonstrations & 122 & 128 \\
    paired subset (same init.\ config) & \multicolumn{2}{c}{102} \\
    \bottomrule
  \end{tabular}
\end{table}

\paragraph{Data.}
We replay the 206 expert demonstrations of the public Push-T dataset~\cite{chi2023diffusion,florence2022implicit} through both domain simulators and keep the successful rollouts: 122 robot-domain and 128 human-domain demonstrations (14.5k/15.5k steps), of which 102 form \emph{paired} episodes sharing the same initial configuration --- the pairs make the mismatch spectrum $\delta_j$ directly estimable from action residuals.
A mirror-convention subtlety is worth recording: the physically correct parity transform maps the block angle as $\theta \mapsto -\theta$ (reflection about the handle axis), not $\theta \mapsto \pi - \theta$; the wrong convention silently destroys replay success (0/206 vs.\ 206/206).
Deviations from~\cite{egobridge2025}: our demonstrations are replayed from public expert data rather than generated by a scripted planner; friction differs across domains in our construction; and we add the paired subset for mismatch estimation.

\paragraph{Policy and training.}
The policy is the CNN-based diffusion policy of~\cite{chi2023diffusion} (2 observation steps, 8 action steps), with a \emph{mixture dataset} that samples human-domain sequences with probability $\lambda$ and robot-domain sequences with $1-\lambda$; validation is always on the robot domain.
Alignment strength is controlled by an optional Gaussian-kernel MMD penalty~\cite{gretton2012kernel} on the shared visual encoder's features, driven by domain labels, with weight $\alpha_{\mathrm{mmd}}$; $\alpha_{\mathrm{mmd}} = 0$ is naive mixing.

\paragraph{Estimating the model's parameters on the testbed.}
The theory's quantities are measured from the dual-domain data: spectra $\hat\alpha_r, \hat\alpha_h$ by PCA on state-action trajectories (mid-range log-log fit); the task-energy spectrum $\hat\beta^{*2}$ by ridge-regressing actions on states in the robot domain; the mismatch spectrum $\hat\delta$ from paired action residuals, yielding $\Delta_\varphi = \|\hat\delta\|^2$ and $c = \hat\delta^\top\hat\beta^*$; effective noise $\tilde\sigma_e^2$ from regression residuals plus feedthrough energy.
On our testbed: $\hat\alpha_r = 1.89$, $\hat\alpha_h = 1.92$ (near-zero spectral gap), and $c < 0$ --- the mirrored construction is indeed \emph{anti-aligned}, so the testbed sits in the regime where the theory predicts only a small marginal benefit, $\lambda^* \approx 0.12$.
